\chapter{CONCLUSION}
\section{Summary of the Present Contribution}
Follow My Voice connects a personalized voice comparison to an instrumented mixed-reality search task. The participant does not know the target identity in advance. Instead, gaze-contingent speech guides orientation and within-wall refinement, while the right controller provides pointing and explicit trigger selection. This organization makes the agent relevant to the search without treating gaze as a selection command.

The system separates spatial evidence, state transitions, constrained model choice, scheduling, and delivered speech. That separation is important when evaluating delayed or contradictory guidance: a model proposal, a qualified gaze event, a controller hold, and a successful selection are different observations. Recording them distinctly makes the interaction more auditable and supports a more precise account of what went wrong when assistance is unhelpful.

The proposed study compares a preferred catalog voice and a self-similar synthetic voice in two counterbalanced 40-trial blocks, following one practice trial in the first assigned voice. Its primary outcome is active completion time adjusted for starting geometry and participant clustering. Retrospective workload, motivation, and agent-experience measures will complement performance after both blocks. The comparison will concern voice identity within a shared adaptive policy, not an isolated benefit of gaze contingency or memory.

At this stage, the contribution is the system, operational design, and measurement plan. Confirmatory participant results are not available, and the thesis makes no claim that either voice improves performance. Before collection, the study must resolve exact pre-cue geometry, balanced vertical allocation, sample-size planning, pilot acceptance criteria, and consent alignment with cloud processing. The completed thesis will then be able to connect its conclusions to actual participant evidence, with uncertainty and limitations stated alongside any observed differences.

\section{Status of the Research Questions}
RQ1 remains unanswered empirically: no confirmatory voice-condition contrast in active completion time is estimated here. RQ2 likewise remains open because the final retrospective workload and motivation responses have not been collected for this report. RQ3 has a defined exploratory battery and debrief structure, but no participant judgments or themes are presented as findings. The literature motivates these questions; it does not supply their answers for this task.

For EQ1, the platform and data contract provide a way to inspect sensing, state, decisions, and software playback separately. The October 20 alignment audit summarized in Chapter~\ref{ch:evaluation} supports specific implementation checks. The separate October 23 working-file inspection in Appendix~\ref{app:decision-contract} documents request, validation, and memory behavior without establishing deployment. It does not establish a held-out guidance error rate or prove that all delayed confirmations have been removed. The contribution is therefore the implemented inspection mechanism and an explicit evaluation design, with its performance claim still bounded by the available record.

\section{Next Empirical Step}
The immediate work is to freeze the study configuration, close the geometry and consent gaps, validate the full wearer procedure, and finalize the primary analysis and sample-size justification before confirmatory recruitment. The resulting data must preserve counterbalanced order, actual voice exposure, observation coverage, and the final retrospective survey references. Only then can this thesis replace its results framework with estimates and revise its abstract and conclusion around observed evidence.

Longer-term experiments could vary memory or provider autonomy independently, as discussed in Chapter~\ref{ch:discussion}. Those would answer different questions from the present voice comparison. Keeping that boundary explicit makes the first study interpretable and provides a concrete basis for subsequent system research.
